Consider two positions, $i$ and $j$, in a protein family.  A structural biologist may draw a line
between them because their atoms lie within a chosen distance.  A statistical physicist may fit a
Potts model and attach a $q\times q$ coupling table to the same pair.  A Transformer may send a
large attention weight from $j$ to $i$.  A model audit may find that substituting the residue at
$j$ changes the predicted distribution at $i$.  A double-mutant experiment may report nonadditive
fitness.  Each result is naturally drawn as an edge.  The pictures look alike, and that visual
agreement invites a stronger conclusion: perhaps all five procedures have found the same
interaction.

They have not---at least not yet.  The structural line is defined by geometry.  The Potts table is
a component of a probability model for an ensemble of homologous sequences.  Attention is a route
inside a computation.  The substitution response is a directed input--output property of a fitted
function.  Experimental epistasis is a statement about an intervention and a measured biological
endpoint.  Agreement among these objects can be scientifically important.  Shared indices $(i,j)$
provide the comparison point; the carried objects and validation criteria determine the meaning.

This book begins with that apparently minor distinction because much of the conceptual difficulty
in protein machine learning is concentrated inside it.  Words such as \emph{interaction},
\emph{energy}, \emph{memory}, \emph{direct}, and \emph{graph} travel easily between statistical
physics and neural computation.  The equations often travel with them, while their assumptions
must be carried explicitly.

\section{A residue pair enters five descriptions}

Let $X=(X_1,\ldots,X_L)$ denote a protein sequence over an alphabet of size $q$.  Suppose that a
multiple-sequence alignment, a trained protein language model, a three-dimensional structure, and
a mutational assay are all available for one family.  For the pair $(i,j)$ they provide five
candidate edge objects:
\begin{figure}[htbp]
\centering
\small
\begin{tabular}{>{\raggedright\arraybackslash}p{0.18\linewidth}
                >{\raggedright\arraybackslash}p{0.27\linewidth}
                >{\raggedright\arraybackslash}p{0.43\linewidth}}
\toprule
Description & Edge object & Question answered \\
\midrule
Structure & contact indicator or geometric relation & Are the two sites close under a declared
structural criterion? \\
Potts model & coupling table $J_{ij}(a,b)$ & Which pair term helps a fitted family distribution
explain sequence variation? \\
Attention & route coefficient $\alpha_{ij}^{(h,\ell)}(x)$ & How much source-normalized weight does
one head assign to this computational route? \\
Model response & categorical map $G_{i\leftarrow j}(C)$ & How does a substitution at $j$ change
the model prediction at $i$ in background $C$? \\
Experiment & epistasis or another intervention response & Does the joint perturbation differ from
the declared additive biological null? \\
\bottomrule
\end{tabular}
\caption{The same endpoints support five different edge semantics.  Agreement must be established
after the carried object and validation question have been declared.}
\label{fig:five-edge-semantics}
\end{figure}

The entries differ along several axes at once.  Some are scalar and some are categorical tables.
Some are symmetric and some directed.  Some are fixed across a family and some depend on the input
sequence.  Some belong to a joint probability law; others belong to an algorithm or an assay.
Even their zeros mean different things.  A zero Potts block may express conditional factorization
within a fitted model, a zero route may mean skipped computation, and a missing structural contact
may only mean that one distance threshold was not crossed.

The first principle of the book is therefore simple:
\begin{equation}
\boxed{\text{shared endpoints do not imply a shared edge semantics.}}
\label{eq:book-shared-endpoints}
\end{equation}
This principle sets the standard for unification: a successful comparison must preserve the
relevant types.  If two edge objects are to be compared, their state spaces, normalization,
conditioning information, invariances, dynamics, and validation criteria must remain visible
during the comparison.

The same caution applies to familiar formulas and familiar adjectives.  Softmax can normalize
residue identities, source positions, memories, or complete configurations.  The normalization
axis determines which probability law the shared algebra defines.  Likewise, an edge may be called \emph{direct}
because it is a one-step computational route, a term in a declared factorization, a spatial
contact, or an experimental effect that survives a lower-order null.  Chapter~2 separates these
meanings systematically.  Here it is enough to retain one reading rule: before interpreting a
formula or an edge, name the space in which it lives and the alternative explanations that have
actually been excluded.

\section{From an edge weight to an edge operator}

For categorical sequences, a scalar edge weight is a compression of a richer object.  If residue
$b$ at site $j$ is replaced by $b'$, the effect at site $i$ is a change across all possible target
identities $a$.  A pair relation is therefore naturally represented by a table or
operator
\begin{equation}
G_{i\leftarrow j}(a,b),
\end{equation}
mapping source substitutions to target contrasts.

This operator contains information that every scalar compression discards.  Two edges can have
the same Frobenius norm while favoring different amino-acid exchanges.  They can have the same
support while acting in different singular directions.  They can generate the same contact score
but make opposite predictions for a particular mutation.  A graph of scalar strengths is useful,
but it is a shadow of the typed relation carried by the edge.

Raw categorical tables introduce another problem.  A constant term, a function of $a$ alone, and
a function of $b$ alone can be moved between one-body and pair terms without changing the modeled
law.  The pair object is therefore an equivalence class until a reference geometry is declared.
The reference-centered operator developed later in the book removes those lower-order directions and
retains the genuinely bivariate action.  Gauge fixing supplies the identified coordinates that
make edge operators comparable across models and references.

\section{One family graph, many sequence responses}

The contrast between a Potts model and a protein language model supplies the running mathematical
problem.  In an additive pairwise Potts model, a fixed gauged block describes the response of site
$i$ to site $j$ throughout the family model.  A nonlinear protein language model generally
produces a different effective response in each sequence background $C$:
\begin{equation}
G_{i\leftarrow j}(C)
=\overline G_{i\leftarrow j}+\delta G_{i\leftarrow j}(C).
\label{eq:central-thesis}
\end{equation}
The mean $\overline G_{i\leftarrow j}$ is the best static approximation in a declared reference
geometry.  The residual $\delta G_{i\leftarrow j}(C)$ measures context dependence.  Under a fixed
product reference, its variance can be identified with higher-order functional components that
contain both audited endpoints.  Under correlated empirical backgrounds, additional cross terms
appear and the interpretation must change.

This decomposition separates two questions that are often merged.  Does the model contain a
stable family-like pair response?  How much of its behavior depends on the surrounding sequence?
A model can answer both strongly, one strongly, or neither.  Calling the entire response a
``learned Potts coupling'' erases the distinction.

Directionality adds another layer.  The response from $j$ to $i$ and the response from $i$ to $j$
are separately measurable.  They define a joint pair energy only if they satisfy a reciprocity
condition after reference transport.  A directed response graph is therefore a valid primary
object.  Symmetrization becomes a later operation, justified after its asymmetry has been measured.

\section{The coordinate system to come}

The rest of the book builds a mathematical coordinate system in which these comparisons can be
made without flattening them.  Each model will be described by the variables carried at its nodes,
the support and type of its edges, the context available to an edge, the represented interaction
order, the updated state, the normalization or conserved quantity, the estimation procedure, and
the external endpoint used for validation.

The construction proceeds from local meaning to global claims.  We first establish a translation
language across physics, probability, and computation.  We then define reference-centered
categorical fibers and the operators between them.  Static and dynamic responses lead to
higher-order decompositions, probability-exact edge geometry, and compatibility conditions for
joint energies.  Potts, Gaussian, Hopfield, RBM, and Transformer models are subsequently placed in
the same coordinates.  Only after the edge object has been declared do we turn to estimation,
spectra, architectural realizability, and biological evidence.

The ambition is to make differences between protein graphs mathematically explicit enough that
genuine agreement becomes meaningful.  The edge between $i$ and $j$ begins a sequence of
questions: an edge in which space, carrying what object, conditioned on what information, invariant
under which changes, and validated against which world?
